% !TEX root = ../DoAn.tex
\documentclass[../DoAn.tex]{subfiles}
\begin{document}

Chương~\ref{chapter:Related_works} đã xác định các nhóm chức năng và yêu cầu của hệ thống. Để hiện thực được các yêu cầu đó, chương này trình bày nền tảng lý thuyết và các công nghệ được lựa chọn, kèm lý do lựa chọn so với các phương án thay thế. Nội dung gồm: nền tảng Android và kiến trúc ứng dụng (Mục~\ref{section:3.1}); công nghệ định vị GNSS (Mục~\ref{section:3.2}); công nghệ hiển thị bản đồ, mô hình ba chiều và thực tế tăng cường (Mục~\ref{section:3.3}); công nghệ Optical Flow và xử lý ảnh (Mục~\ref{section:3.4}); backend xử lý Optical Flow bằng mô hình mạng lưới học sâu (Mục~\ref{section:3.5}); cơ chế kết hợp GNSS, Optical Flow và IMU trong LiveRouting (Mục~\ref{section:3.6}); và quy trình phát triển phần mềm (Mục~\ref{section:3.7}).

\section{Nền tảng Android và kiến trúc ứng dụng}
\label{section:3.1}
\subsection{Ngôn ngữ lập trình Kotlin}
\label{subsection:3.1.1}
Ứng dụng được phát triển bằng Kotlin trên nền Android. Kotlin được chọn vì đây là ngôn ngữ chính thức và hiện đại của hệ sinh thái Android, có cú pháp ngắn gọn, hỗ trợ null-safety và phù hợp với lập trình bất đồng bộ bằng coroutine. So với Java, Kotlin giúp giảm mã lặp trong các lớp UI, data class và callback. Điều này đặc biệt hữu ích với đồ án vì hệ thống có nhiều màn hình và nhiều luồng dữ liệu cảm biến chạy song song như vị trí, vệ tinh, camera và IMU.

Android SDK cung cấp các API nền tảng như \texttt{LocationManager}, \texttt{GnssStatus}, \texttt{GnssMeasurementsEvent}, \texttt{SensorManager}, hệ thống quyền, foreground service và lưu trữ. Đây là các thành phần bắt buộc để ứng dụng truy cập vị trí, vệ tinh, camera, cảm biến và dịch vụ chạy nền. Đồ án đã được cấu hình để cân bằng giữa khả năng tương thích với nhiều thiết bị và việc sử dụng các API mới.

\subsection{Kiến trúc giao diện XML Modern Android}
\label{subsection:3.1.2}
Để xây dựng giao diện ổn định, dễ mở rộng và phù hợp với nhiều thành phần đặc thù của hệ thống, ứng dụng sử dụng hướng tiếp cận Modern Android kết hợp layout XML và các thành phần Jetpack.

Layout XML được sử dụng để xây dựng giao diện người dùng theo mô hình View truyền thống. Cách tiếp cận này phù hợp với phạm vi đồ án vì nhiều thành phần quan trọng như \texttt{GLSurfaceView} dùng cho mô hình Trái Đất 3D, preview camera và bản đồ osmdroid đều tích hợp tự nhiên với hệ thống View của Android. Nhờ đó, ứng dụng có thể kết hợp linh hoạt giữa giao diện thông thường, camera thời gian thực, bản đồ và đồ họa OpenGL trong cùng một hệ thống.

Jetpack Navigation được sử dụng để điều phối quá trình chuyển màn hình trong ứng dụng. Hệ thống sử dụng một \texttt{MainActivity} chứa \texttt{NavHostFragment}, trong đó các chức năng chính được chia thành nhiều fragment như màn hình giới thiệu, màn hình chính, GNSS Viewer, LiveRouting, GNSS AR, Camera Optical Flow, danh sách media, trình phát video, danh sách analytics và chi tiết analytics. Việc sử dụng navigation graph giúp luồng điều hướng rõ ràng hơn, đồng thời tránh để Activity phải trực tiếp quản lý toàn bộ fragment.

ViewModel được sử dụng để lưu trữ và quản lý trạng thái giao diện theo vòng đời ứng dụng. ViewModel của GNSS Viewer quản lý chế độ hiển thị 2D/3D, vị trí hiện tại, địa điểm đã chọn, tuyến đường và dữ liệu vệ tinh mới nhất. ViewModel của LiveRouting lưu trạng thái tuyến đường, dữ liệu GNSS, Optical Flow assist, ước lượng di chuyển và kết quả hiệu chỉnh theo tuyến. ViewModel của Camera Optical Flow quản lý thuật toán đang chọn, vùng quan tâm, trạng thái ghi hình, phiên phân tích và chế độ chuyển động. Nhờ đó, ứng dụng hạn chế mất trạng thái khi màn hình được tạo lại, đồng thời tách rõ phần hiển thị khỏi phần xử lý nghiệp vụ.

\section{Công nghệ định vị GNSS}
\label{section:3.2}
\subsection{GNSS trên thiết bị Android}
\label{subsection:3.2.1}
GNSS là nhóm hệ thống vệ tinh dẫn đường toàn cầu, bao gồm GPS, Galileo, BeiDou, GLONASS và một số hệ thống khác. Trên Android, \texttt{GnssStatus} cung cấp thông tin quan sát như SVID, chòm vệ tinh, azimuth, elevation, C/N0 và trạng thái used-in-fix. \texttt{GnssMeasurementsEvent} cung cấp dữ liệu đo thô hơn từ GNSS engine~\cite{android-gnss-event}, trong đó một số thiết bị có thể cung cấp thông tin PVT của vệ tinh thông qua các API bị ẩn hoặc khác biệt theo phiên bản~\cite{android-gnss-measurement}.

Trong đồ án, GNSS không chỉ dùng để lấy vị trí người dùng mà còn dùng để tạo danh sách vệ tinh và dữ liệu hiển thị trên renderer. Mỗi vệ tinh được định danh bằng cặp \texttt{constellationType} và \texttt{svid}. Cặp này tạo thành \texttt{SatelliteKey}, giúp cache PVT, broadcast ephemeris và CelesTrak orbit theo cùng một khóa logic.

\subsection{Hệ tọa độ WGS84, ECEF và LLA}
\label{subsection:3.2.2}
Đồ án sử dụng hệ quy chiếu WGS84~\cite{wgs84-tr8350} cho biểu diễn vĩ độ, kinh độ và độ cao. Khi dữ liệu vệ tinh ở hệ ECEF, ứng dụng chuyển sang LLA bằng công thức ellipsoid WGS84. Với bán trục lớn $a$, bán trục nhỏ $b$, bình phương độ lệch tâm thứ nhất $e^2$ và thứ hai $e'^2$, phép chuyển đổi Bowring được viết rút gọn như Phương trình~\ref{eq:ecef-lla}. Công thức này là nền tảng để hiển thị vị trí vệ tinh trên bản đồ và mô hình 3D.

\begin{equation}
\label{eq:ecef-lla}
\begin{split}
p &= \sqrt{X^2 + Y^2},\quad
\theta = \arctan\frac{Z a}{p b},\\
\varphi &= \arctan\frac{Z + e'^2 b \sin^3\theta}{p - e^2 a \cos^3\theta},\quad
\lambda = \arctan\frac{Y}{X}.
\end{split}
\end{equation}

Trong đó $(X,Y,Z)$ là tọa độ ECEF, $\varphi$ là vĩ độ và $\lambda$ là kinh độ. Sau khi có vĩ độ và kinh độ, độ cao được tính từ bán kính cong pháp tuyến tại vĩ độ tương ứng. Trong mã nguồn, logic này được đặt trong \texttt{SatelliteCalculator.kt}.

\subsection{Dữ liệu quỹ đạo vệ tinh: IGS, CelesTrak và SGP4}
\label{subsection:3.2.3}
Nguồn dữ liệu vệ tinh được tổ chức theo chuỗi ưu tiên. Nguồn thứ nhất là PVT thật từ thiết bị, nếu trích xuất được. Nguồn thứ hai là broadcast ephemeris từ IGS~\cite{igs-products}, một tổ chức cung cấp dữ liệu và sản phẩm GNSS công khai. Nguồn thứ ba là dữ liệu TLE từ CelesTrak GP~\cite{celestrak-gp-data,ccsds-omm}, trong đó ứng dụng tải các nhóm GPS, Galileo và BeiDou rồi ánh xạ tên vệ tinh về SVID Android. TLE phù hợp với bài toán trực quan hóa vì có thể lan truyền bằng SGP4 và kích thước dữ liệu nhỏ.

SGP4 là mô hình lan truyền quỹ đạo từ phần tử trung bình. Trong hệ thống, \texttt{Sgp4OrbitPropagator} được dùng khi có dữ liệu CelesTrak phù hợp. Nếu không thể dùng SGP4, ứng dụng fallback về phương pháp tính từ phần tử Kepler cơ bản. Cách tổ chức này giúp module GNSS không phụ thuộc vào một nguồn dữ liệu duy nhất.

Các phép tính quỹ đạo cần xử lý đúng hệ thời gian. Ứng dụng khởi tạo Orekit bằng cách chép file \texttt{tai-utc.dat} từ assets vào thư mục cache và đăng ký với \texttt{DataProvidersManager}. File này cung cấp dữ liệu chênh lệch TAI-UTC/leap second, giúp \texttt{AbsoluteDate} và \texttt{TimeScalesFactory.getUTC()} làm việc ổn định khi lan truyền broadcast ephemeris hoặc TLE theo thời điểm quan sát.

\section{Công nghệ hiển thị bản đồ, 3D và thực tế tăng cường}
\label{section:3.3}
\subsection{Hiển thị vị trí và vệ tinh trên bản đồ}
\label{subsection:3.3.1}
Bản đồ hai chiều được dựng bằng thư viện osmdroid trên nền dữ liệu OpenStreetMap. osmdroid được chọn vì là thư viện mã nguồn mở, không yêu cầu khóa API thương mại và cho phép tùy biến trực tiếp các lớp marker, overlay và polyline. Trên bản đồ, ứng dụng hiển thị marker vị trí người dùng, kết quả tìm kiếm địa điểm và tuyến đường vẽ từ hình học route. Dịch vụ Nominatim được dùng cho tìm kiếm địa điểm, còn OSRM cung cấp hình học tuyến đường giữa hai điểm.

Ngoài vị trí người dùng, bản đồ còn là một trong các kênh trực quan hóa vệ tinh. Từ azimuth và elevation trong \texttt{GnssStatus}, ứng dụng có thể biểu diễn hướng và độ cao của vệ tinh so với người quan sát, giúp đối chiếu nhanh với mô hình 3D và màn hình AR. Khi vị trí vệ tinh được phân giải từ PVT, IGS hoặc CelesTrak (xem Mục~\ref{subsection:3.2.3}), thông tin nguồn dữ liệu cũng được gắn kèm để người dùng biết độ tin cậy tương đối của từng điểm hiển thị.

\subsection{OpenGL ES trong mô phỏng không gian 3D}
\label{subsection:3.3.2}
Mô hình Trái Đất ba chiều được dựng bằng OpenGL ES. Hệ tọa độ của renderer đặt gốc tại tâm Trái Đất, trục $Y$ hướng Bắc, trục $Z$ trùng kinh tuyến gốc và trục $X$ hướng kinh độ $90^\circ$ Đông. Với vĩ độ $\varphi$, kinh độ $\lambda$ và bán kính hiển thị $r$, tọa độ cầu được đưa lên hệ Descartes theo Phương trình~\ref{eq:spherical-cartesian}.

\begin{equation}
\label{eq:spherical-cartesian}
x = r \cos\varphi \sin\lambda,\quad
y = r \sin\varphi,\quad
z = r \cos\varphi \cos\lambda.
\end{equation}

OpenGL ES được chọn vì cho phép vẽ mô hình 3D hiệu quả trên thiết bị Android, kiểm soát camera ảo, phép quay, zoom và tương tác chạm. Trong \texttt{EarthRenderer}, Trái Đất dùng texture ban ngày và ban đêm, shader tính hướng chiếu sáng từ Mặt Trời để tạo vùng sáng/tối và đường chuyển tiếp ngày-đêm. Renderer cũng tính vị trí tương đối của Mặt Trời và Mặt Trăng theo UTC để vẽ các thiên thể này trong cùng cảnh 3D. Phương án thay thế là sử dụng một engine 3D hoàn chỉnh như Unity hoặc Sceneform. Tuy nhiên, các engine này làm tăng kích thước dự án và phụ thuộc runtime. Với phạm vi đồ án, OpenGL ES trực tiếp đủ để dựng Trái Đất, vệ tinh, Mặt Trời, Mặt Trăng và các tương tác cơ bản.

\subsection{Thực tế tăng cường trong quan sát vệ tinh}
\label{subsection:3.3.3}
Ở màn hình AR, ứng dụng dùng ARCore để lấy frame camera, pose và projection matrix, đồng thời dùng cảm biến hướng thiết bị để hiệu chỉnh yaw theo hướng nhìn của camera sau. Cách tiếp cận này không nhằm tạo một hệ AR định vị chính xác tuyệt đối như các SDK chuyên sâu, mà nhằm minh họa vị trí tương đối của vệ tinh trên hướng nhìn của người dùng. Việc giữ module AR tách khỏi module bản đồ giúp hệ thống có thể dùng chung dữ liệu vệ tinh nhưng thay đổi cách biểu diễn, từ bản đồ phẳng sang mô hình 3D và sang khung nhìn AR.

\section{Công nghệ Optical Flow và xử lý ảnh}
\label{section:3.4}
\subsection{Cơ sở lý thuyết Optical Flow}
\label{subsection:3.4.1}
Optical Flow ước lượng chuyển động biểu kiến của điểm ảnh giữa hai khung hình liên tiếp. Giả thiết cơ bản là độ sáng của một điểm không đổi trong khoảng thời gian ngắn, được mô tả bởi $I(x,y,t)=I(x+\Delta x,y+\Delta y,t+\Delta t)$. Khai triển Taylor bậc nhất dẫn tới ràng buộc gradient trong Phương trình~\ref{eq:brightness-constancy}.

\begin{equation}
\label{eq:brightness-constancy}
I_x u + I_y v + I_t = 0.
\end{equation}

Phương trình trên có hai ẩn $u$ và $v$ nhưng chỉ có một ràng buộc tại mỗi điểm ảnh, nên cần thêm giả thiết hoặc mô hình cục bộ để giải. Các thuật toán Optical Flow khác nhau chủ yếu ở cách bổ sung giả thiết này, cách lấy mẫu điểm ảnh và cách tối ưu kết quả.

\subsection{Thư viện OpenCV}
\label{subsection:3.4.2}
OpenCV là thư viện thị giác máy tính mã nguồn mở được dùng làm nền tảng xử lý ảnh trong đồ án. Trên Android, OpenCV cung cấp kiểu dữ liệu \texttt{Mat} để biểu diễn ảnh và một tập hàm Optical Flow sẵn có như \texttt{goodFeaturesToTrack}, \texttt{calcOpticalFlowPyrLK} và \texttt{calcOpticalFlowFarneback}~\cite{opencv-optical-flow}. Nhờ vậy, ứng dụng có thể tập trung vào việc thiết kế luồng xử lý và hiển thị thay vì hiện thực lại các thuật toán nền tảng.

OpenCV được chọn vì có hiệu năng tốt, được kiểm chứng rộng rãi và có cùng bộ thuật toán trên cả thiết bị Android lẫn máy chủ Python, giúp giữ tính nhất quán giữa Optical Flow chạy trên thiết bị và pipeline xử lý video phía server. Trong ứng dụng, mỗi frame camera được chuyển sang \texttt{Mat}, xử lý bằng KLT hoặc Farneback rồi phủ kết quả lên khung hình hiển thị.

\subsection{Thuật toán KLT và Farneback}
\label{subsection:3.4.3}
KLT trong đồ án kết hợp phát hiện đặc trưng Shi-Tomasi~\cite{shi-tomasi1994} và theo dõi Lucas-Kanade~\cite{lucas-kanade1981} dạng kim tự tháp. Shi-Tomasi chọn các điểm có cấu trúc gradient đủ tốt để theo dõi, còn Lucas-Kanade giả định chuyển động cục bộ gần như không đổi trong một cửa sổ nhỏ và giải hệ bình phương tối thiểu. Cách tiếp cận này xử lý ít điểm nên phù hợp với yêu cầu thời gian thực trên thiết bị di động. Trong mã nguồn, lớp \texttt{KLT} chuyển frame sang grayscale, phát hiện điểm bằng \texttt{goodFeaturesToTrack}, theo dõi bằng \texttt{calcOpticalFlowPyrLK} và lọc kết quả bằng kiểm tra tiến-lùi. Các vector quá nhỏ bị xem là nhiễu và không được đưa vào motion vector. Kết quả dịch chuyển được làm mượt bằng exponential smoothing để giảm rung khi hiển thị.

Farneback~\cite{farneback2003} thuộc nhóm dense optical flow. Thuật toán xấp xỉ vùng lân cận của ảnh bằng đa thức bậc hai và suy ra dịch chuyển từ sự biến đổi của các hệ số đa thức giữa hai khung hình. Trong đồ án, Farneback được chạy trên ảnh giảm kích thước để giảm chi phí xử lý, sau đó vector được lấy mẫu theo lưới để vẽ lên khung hình. Chế độ heatmap dùng độ lớn vector để tạo bản đồ màu, phù hợp khi người dùng muốn quan sát phân bố chuyển động thay vì từng mũi tên riêng lẻ. Farneback có chi phí cao hơn KLT vì phải tính trường chuyển động dày đặc, đổi lại không phụ thuộc hoàn toàn vào một số điểm đặc trưng rời rạc. Vì vậy, hai thuật toán bổ trợ cho nhau: KLT phù hợp môi trường nhiều đặc trưng rõ và cần tốc độ, còn Farneback phù hợp để minh họa chuyển động tổng thể của ảnh. Trong đồ án, người dùng có thể chuyển giữa vector và heatmap để quan sát hai kiểu biểu diễn khác nhau.

\subsection{Phân tích chuyển động từ camera thời gian thực}
\label{subsection:3.4.4}
Luồng camera thời gian thực được cung cấp bởi CameraX, thư viện gắn trực tiếp với lifecycle của Android và cung cấp use case \texttt{ImageAnalysis}~\cite{android-camerax-imageanalysis} cho xử lý ảnh truy cập được trên CPU. Trong đồ án, CameraX tạo hai use case: \texttt{Preview} để hiển thị camera gốc và \texttt{ImageAnalysis} để đưa frame vào OpenCV. Chiến lược backpressure \texttt{STRATEGY\_KEEP\_ONLY\_LATEST} được sử dụng để tránh tích tụ frame khi thuật toán xử lý chậm. Phương án thay thế là dùng Camera2 API trực tiếp; tuy linh hoạt hơn nhưng Camera2 phức tạp đáng kể khi quản lý session, surface và lifecycle, nên CameraX phù hợp hơn với mục tiêu tập trung vào thuật toán của đồ án.

Để đánh giá độ tin cậy tương đối của vector trong thời gian thực, đồ án sử dụng Forward-Backward Error. Thuật toán chạy flow từ frame trước sang frame sau, sau đó chạy ngược từ frame sau về frame trước. Nếu điểm quay về cách điểm ban đầu dưới ngưỡng $1.5$ pixel, vector được xem là inlier. Độ tin cậy được tính theo Phương trình~\ref{eq:fbe-confidence}.

\begin{equation}
\label{eq:fbe-confidence}
\text{Confidence} = \frac{N_{\text{inlier}}}{N_{\text{tracked}}} \times 100\%.
\end{equation}

Chỉ số này không phải là độ chính xác tuyệt đối so với ground truth, vì hệ thống không có dữ liệu ground truth của chuyển động camera. Tuy nhiên, nó giúp phát hiện các vector thiếu nhất quán và tạo cơ sở so sánh tương đối giữa KLT và Farneback trong cùng phiên chạy, phục vụ cho cơ chế phân tích hiệu năng được trình bày ở Chương~\ref{chapter:SolutionAndContribution}.

\section{Backend xử lý Optical Flow bằng mô hình mạng lưới học sâu}
\label{section:3.5}
\subsection{FastAPI trong xây dựng server xử lý video}
\label{subsection:3.5.1}
FastAPI được chọn cho máy chủ Python tự triển khai vì hỗ trợ tốt multipart upload~\cite{fastapi-upload}, background task~\cite{fastapi-background} và kiểu dữ liệu rõ ràng cho API. Trong đồ án, endpoint tạo job được dùng làm luồng chính cho xử lý bất đồng bộ; các endpoint kiểm thử chỉ phục vụ phát triển server. Cách tổ chức này tránh việc thiết bị Android phải giữ một kết nối HTTP dài trong toàn bộ thời gian xử lý video. Máy chủ được chạy bằng Uvicorn, một ASGI server nhẹ, phù hợp cho cả môi trường phát triển và demo.

\subsection{Mô hình RAFT và ONNX Runtime}
\label{subsection:3.5.2}
RAFT~\cite{raft2020} là mô hình học sâu cho Optical Flow, xây dựng volume tương quan toàn cặp giữa hai ảnh và cập nhật trường flow lặp bằng một đơn vị recurrent. So với KLT và Farneback, RAFT có khả năng biểu diễn chuyển động dày đặc và phức tạp hơn, nhưng yêu cầu suy luận nặng hơn. Vì vậy, đồ án triển khai RAFT theo hướng backend server cho xử lý video ngoại tuyến.

Mô hình RAFT được đóng gói ở dạng ONNX. ONNX Runtime~\cite{onnxruntime-python-api} cung cấp lớp \texttt{InferenceSession} để nạp mô hình và chạy suy luận, đồng thời hỗ trợ chọn execution provider như CPU, CUDA hoặc DirectML khi môi trường có sẵn. Máy chủ \texttt{OpticalFlowServer} dùng ONNX Runtime thay vì OpenCV DNN cho pipeline Python vì cần xử lý mô hình ONNX quantized ổn định hơn trong bối cảnh server.

\subsection{Thư viện Cutie trong xử lý ROI đối tượng}
\label{subsection:3.5.3}
Khi người dùng chọn vùng quan tâm (ROI) trên video, backend mặc định dùng Cutie để lan truyền mask đối tượng. Cutie là mô hình phân vùng và theo dõi đối tượng trong video: từ một mask khởi tạo tại frame được chọn, mô hình lan truyền mask qua các frame khác để bám đối tượng ngay cả khi nó dịch chuyển hoặc biến dạng. Trong pipeline, Optical Flow sau đó chỉ được vẽ trong vùng mask đang hoạt động, giúp tập trung vào đối tượng và giảm nhiễu nền. Nếu cần triển khai đơn giản hơn hoặc môi trường không đủ tài nguyên, backend có thể chuyển về chế độ ROI tĩnh bằng biến môi trường.

\subsection{Cơ chế xử lý video theo job và chunk}
\label{subsection:3.5.4}
Vì video là file lớn và RAFT có thời gian suy luận cao, việc gửi video qua một request đồng bộ duy nhất dễ gây timeout và không cho biết tiến độ. Do đó, backend được thiết kế theo cơ chế job bất đồng bộ: mỗi yêu cầu xử lý sinh ra một \texttt{job\_id}, video được đưa vào hàng xử lý nền và client thăm dò trạng thái theo chu kỳ. Với file lớn, ứng dụng tạo upload session và gửi video theo từng chunk (kèm kiểm tra toàn vẹn rồi ráp lại ở máy chủ); kết quả cũng có thể tải về theo chunk khi file đầu ra lớn.

Phía Android, Retrofit~\cite{retrofit-docs} được dùng cho giao tiếp HTTP với backend vì cho phép khai báo API bằng interface, hỗ trợ multipart request và streaming response, đồng thời kết hợp tốt với coroutine. \texttt{OpticalFlowServerApi} định nghĩa các endpoint tạo job, tạo upload session, gửi chunk, hoàn tất upload, lấy trạng thái job, hủy job, lấy thông tin kết quả, tải kết quả và dọn kết quả. Tác vụ upload, polling và download được thực thi bởi một foreground worker dựa trên WorkManager, nhờ đó tiến trình không bị dừng khi người dùng rời màn hình. Ở phía server, kết quả được ghi ra MP4 H.264 bằng \texttt{imageio-ffmpeg} hoặc ffmpeg trong PATH để Android phát ổn định, thay cho việc phụ thuộc vào \texttt{VideoWriter} của OpenCV cho container đầu ra. Khi tải về, video được mở bằng Media3/ExoPlayer, vốn xử lý tốt các định dạng media hiện đại và cho phép tùy biến giao diện phát video trong fragment.

\section{Kết hợp GNSS, Optical Flow và IMU trong LiveRouting}
\label{section:3.6}
\subsection{Cảm biến IMU trên thiết bị Android}
\label{subsection:3.6.1}
IMU (Inertial Measurement Unit) trên điện thoại gồm gia tốc kế và con quay hồi chuyển, được truy cập qua \texttt{SensorManager} của Android. Trong đồ án, dữ liệu IMU phục vụ hai mục đích chính. Thứ nhất, tốc độ quay quanh trục đứng (yaw rate) từ con quay hồi chuyển được dùng để ước lượng thay đổi hướng của thiết bị khi GNSS không cập nhật kịp. Thứ hai, gia tốc kế giúp nhận biết trạng thái tĩnh hay đang di chuyển; khi phát hiện thiết bị gần như đứng yên, hệ thống áp dụng cơ chế ZUPT (Zero-velocity Update) để khử trôi vận tốc tích lũy. IMU có ưu điểm là tần số cao và không phụ thuộc tín hiệu ngoài, nhưng dễ tích lũy sai số theo thời gian, nên chỉ phù hợp để hỗ trợ ngắn hạn và cần được ràng buộc bởi GNSS hoặc Optical Flow.

\subsection{Visual Odometry thực dụng}
\label{subsection:3.6.2}
Visual odometry là bài toán ước lượng chuyển động tương đối của camera từ chuỗi ảnh liên tiếp. Trong dạng đầy đủ, hệ thống thường cần theo dõi đặc trưng, loại outlier, ước lượng pose camera, xử lý tỉ lệ mét/pixel và kết hợp IMU hoặc bản đồ để giảm trôi. Tuy nhiên, với camera đơn trên điện thoại và không có hiệu chuẩn độ sâu hoặc chiều cao camera ổn định, Optical Flow chỉ cung cấp chuyển động biểu kiến trong ảnh; nó không trực tiếp tương đương với vận tốc thật theo mét/giây.

Trong đồ án, LiveRouting dùng một phiên bản visual odometry thực dụng cho kịch bản xe. Metrics Optical Flow gồm \texttt{avgDx}, \texttt{avgDy}, \texttt{avgMagnitude} và \texttt{confidence} được quy đổi thành tốc độ tương đối, thay đổi hướng và hệ số chất lượng. Tỉ lệ chuyển đổi pixel/giây sang mét/giây được hiệu chỉnh khi GNSS còn tốt và xe có tốc độ đủ lớn. Khi GNSS mất tin cậy, tốc độ ước lượng được lấy từ Optical Flow với trọng số phụ thuộc chất lượng; tốc độ GNSS cuối cùng chỉ còn là prior và giảm dần theo thời gian để tránh giả định xe luôn giữ nguyên vận tốc.

\subsection{Tùy chọn bám tuyến khi mất GNSS}
\label{subsection:3.6.3}
Khi GNSS suy giảm, vị trí được suy ra hoàn toàn từ dead reckoning (Optical Flow, IMU và con quay hồi chuyển). Để giảm sai số hiển thị, đồ án cung cấp một tùy chọn bám tuyến (chế độ SNAP): khi vị trí ước lượng nằm đủ gần tuyến đã vạch, marker được kéo nhẹ về tuyến và tiến theo tuyến tương ứng với quãng đường đã đi; khi độ lệch vượt ngưỡng, hệ thống nhả khóa để marker đi theo dự đoán. Đây không phải map matching khớp liên tục vào mạng đường, mà chỉ là cơ chế bám vào một tuyến đã chọn và có thể tắt bằng chế độ REAL để hiển thị thẳng vị trí dead reckoning nhằm thấy rõ độ lệch. Nhờ tách rời như vậy, hệ thống vẫn giữ được bằng chứng lệch tuyến để kích hoạt cập nhật lại đường đi qua OSRM khi người dùng thật sự đi sai đường.

\subsection{Cơ chế hỗ trợ định vị khi GNSS suy giảm}
\label{subsection:3.6.4}
Ba nguồn dữ liệu trên được kết hợp thành một cơ chế hỗ trợ định vị ngắn hạn cho LiveRouting. Khi GNSS còn tin cậy (đủ số vệ tinh used-in-fix, độ chính xác tốt và dữ liệu mới), vị trí và tốc độ từ Android Location API là nguồn chính, đồng thời hệ thống tranh thủ hiệu chỉnh hệ số chuyển đổi pixel/giây sang mét/giây cho Optical Flow và trục tiến của IMU. Khi GNSS suy giảm, hệ thống chuyển sang ước lượng bằng dead reckoning: tốc độ lấy từ visual odometry kết hợp gia tốc kế IMU, hướng được cập nhật bằng yaw rate của con quay hồi chuyển, và vị trí dự đoán được giới hạn bởi các ngưỡng tăng tốc, phanh và coast để tránh biến động phi thực tế. Vị trí dự đoán sau đó có thể được bám vào tuyến qua tùy chọn bám tuyến (chế độ SNAP) khi nằm đủ gần tuyến, hoặc giữ nguyên ở chế độ REAL khi nghi ngờ lệch tuyến. Sự kết hợp này nhằm duy trì marker liên tục trong các khoảng mất GNSS ngắn, nhưng vẫn được trình bày đúng bản chất là một cơ chế hỗ trợ thực dụng, không phải hệ visual-inertial odometry đầy đủ. Chi tiết thuật toán fusion được trình bày ở Chương~\ref{chapter:SolutionAndContribution}.

\section{Quy trình phát triển phần mềm}
\label{section:3.7}
Đồ án được phát triển theo hướng linh hoạt (Agile) với các vòng lặp ngắn phù hợp quy mô một sinh viên thực hiện. Các hạng mục chức năng (quan sát vệ tinh, chỉ đường trực tiếp, camera Optical Flow, phân tích hiệu năng và xử lý video bằng mô hình mạng lưới học sâu) được làm theo thứ tự ưu tiên và phụ thuộc kỹ thuật, trong đó GNSS và Optical Flow thời gian thực được hoàn thiện trước vì là nền tảng cho các module còn lại. Mỗi vòng lặp gồm phân tích, thiết kế, hiện thực, kiểm thử trên thiết bị thật rồi tinh chỉnh theo kết quả quan sát. Mã nguồn được quản lý bằng Git, ứng dụng build bằng Android Studio và Gradle, còn máy chủ Python chạy thử cục bộ bằng Uvicorn. Do hệ thống phụ thuộc nhiều vào phần cứng và cảm biến, việc kiểm thử chủ yếu là kiểm thử chức năng và tích hợp trên thiết bị thật kết hợp quan sát log (chi tiết ở Chương~\ref{chapter:Experiment}).

\section*{Kết chương}
Chương này đã trình bày các nền tảng lý thuyết và công nghệ của đồ án: ngôn ngữ Kotlin và kiến trúc Modern Android cho ứng dụng; chuỗi nguồn dữ liệu GNSS (PVT, IGS, CelesTrak/SGP4) cùng các phép biến đổi tọa độ; osmdroid, OpenGL ES và ARCore cho ba kênh trực quan hóa vệ tinh; KLT, Farneback và RAFT cho Optical Flow; FastAPI, ONNX Runtime và Cutie cho backend xử lý video; cùng cơ chế kết hợp GNSS, Optical Flow và IMU theo hướng dead reckoning cho LiveRouting. Những lựa chọn công nghệ này là tiền đề trực tiếp cho phần thiết kế, xây dựng, kiểm thử và triển khai hệ thống, được trình bày trong Chương~\ref{chapter:Experiment}.

\end{document}
